\section*{Capítulo \ref{ch:b1:rifts-ridges-margins} --- Riftes, dorsais e margens passivas}

\begin{solution}{exo:b1:rifts-ridges-margins:1}
(a) O dique cobriria $\qty{60000}{m}\times\qty{8}{m} = \qty{4.8e5}{m^2}$ em
planta, logo sua altura média seria $\qty{2.5e9}{m^3}/\qty{4.8e5}{m^2} =
\qty{5200}{m}$, cerca de \qty{5}{km}: o dique cortou grande parte da \omterm{def:b1:layered-earth:layers}{crosta}.
(b) $\qty{8}{m}/\qty{0.018}{m.yr^{-1}} = 444$ anos, uns quatro séculos.
\end{solution}

\begin{solution}{exo:b1:rifts-ridges-margins:2}
(a) O \omterm{def:b1:cross-sections:throw}{afastamento horizontal} é o \omterm{def:b1:cross-sections:throw}{rejeito} vertical dividido pela tangente do \omterm{def:b1:geological-maps:attitude}{mergulho}:
$3.0/\tan\ang{60} = \qty{1.73}{km}$ e $0.8/\tan\ang{70} = \qty{0.29}{km}$;
a extensão da seção é \qty{2.02}{km}. (b) Antes do falhamento a seção tinha
$25 - 2.02 = \qty{22.98}{km}$ de largura, cerca de \qty{23}{km}: foi alongada em
cerca de \qty{9}{\percent}.
\end{solution}

\begin{solution}{exo:b1:rifts-ridges-margins:3}
Ultralenta (abaixo de \qty{20}{mm.yr^{-1}}): a Dorsal de Gakkel (12), o centro do
mar Vermelho (15) e a Dorsal de Reykjanes (19), esta última perto do limite. Lenta (20 a
\qty{55}{mm.yr^{-1}}): a Dorsal de Carlsberg ocidental (23) e o Atlântico Sul
(33). Rápida (acima de \qty{75}{mm.yr^{-1}}): a Dorsal do Pacífico Leste (143). As
dorsais ultralentas e lentas devem ter um vale axial entre falhas normais; a Dorsal do
Pacífico Leste, um alto axial.
\end{solution}

\begin{solution}{exo:b1:rifts-ridges-margins:4}
(a) $\qty{3.0}{km^2}\times\qty{7.1}{km} = \qty{21.3}{km^3}$, cerca de
\qty{21}{km^3} por ano. (b) $\qty{3.028e8}{km^2}/\qty{3.0}{km^2.yr^{-1}} =
\qty{1.01e8}{yr}$, \qty{101}{Myr}. O fundo mais antigo, cerca de \qty{190}{Myr}, é
mais antigo porque o fundo não é destruído em ordem de idade: a subducção consome
fundo de todas as idades, parte do fundo antigo sobrevive onde nenhuma fossa o alcançou,
e a taxa de produção variou.
\end{solution}

\begin{solution}{exo:b1:rifts-ridges-margins:5}
(a) $2500 + 350\times2 = \qty{3200}{m}$, $2500 + 350\times6 = \qty{4600}{m}$,
$2500 + 350\times8 = \qty{5300}{m}$. (b) $t = \bigl((d - 2500)/350\bigr)^2$:
$(2100/350)^2 = \qty{36}{Myr}$ e $(2800/350)^2 = \qty{64}{Myr}$. (c) O
fundo afunda \qty{1400}{m} entre 4 e \qty{36}{Myr}, mas apenas \qty{700}{m}
entre 36 e \qty{64}{Myr}, em um intervalo de tempo semelhante.
\end{solution}

\begin{solution}{exo:b1:rifts-ridges-margins:6}
O coeficiente de expansão efetivo é
$\num{3.2e-5}\times3300/2300 = \qty{4.591e-5}{K^{-1}}$, e
$\sqrt{\kappa\times\qty{1}{Myr}} = \sqrt{\num{e-6}\times\num{3.156e13}} =
\qty{5618}{m}$. Então $C = (2/\sqrt\pi)\times\num{4.591e-5}\times1350\times5618 =
1.1284\times\num{4.591e-5}\times1350\times5618 = \qty{393}{m.Myr^{-1/2}}$. Para um
\omterm{def:b1:layered-earth:layers}{manto} \qty{250}{K} mais frio, $C$ varia com $T_m - T_0$:
$393\times1100/1350 = \qty{320}{m.Myr^{-1/2}}$. O fundo de um oceano mais frio
afunda mais devagar e segue a constante menor. O \num{350} ajustado do
capítulo não é um desses casos: ele está dentro da incerteza dos parâmetros.
\end{solution}

\begin{solution}{exo:b1:rifts-ridges-margins:7}
(a) $\qty{e-6}{cal.cm^{-2}.s^{-1}} = \num{4.184e-6}\times\num{e4} =
\qty{4.184e-2}{W.m^{-2}} = \qty{41.84}{mW.m^{-2}}$, logo a constante é
$11.3\times41.84 = \qty{473}{mW.m^{-2}.Myr^{1/2}}$, \qty{3.5}{\percent} abaixo de
os 490 de \cref{ch:b1:heat-budget}: os dois ajustes concordam. (b) $q = 473$,
$473/\sqrt{10} = 150$ e $473/\sqrt{150} = \qty{38.6}{mW.m^{-2}}$. A lei
da raiz quadrada continua caindo, enquanto o modelo de placa se nivela em uma constante:
a \qty{150}{Myr} o valor pela raiz quadrada é baixo demais, pois a lei de profundidade vai fundo demais.
\end{solution}

\begin{solution}{exo:b1:rifts-ridges-margins:8}
Tome o segmento norte sobre $x = 0$ e o sul sobre $x = -40$ (km), a transformante sobre $y = 0$ entre eles. (a) Logo ao norte da transformante o fundo
fica a oeste do segmento norte e se move para oeste; logo ao sul dela, fica a leste do segmento sul e se move para leste. Vistos de qualquer dos lados, o lado distante se move para a esquerda: o deslizamento é sinistral (levógiro). (b) A dorsal parece ter sido deslocada dextralmente, sua parte sul a oeste: o sentido oposto, como prevê \cref{prop:b1:rifts-ridges-margins:transform}.
(c) Terremotos ocorrem apenas nos \qty{40}{km} entre as duas extremidades da dorsal, onde os dois lados pertencem a placas diferentes; além, ao longo das \omterm{def:b1:rifts-ridges-margins:transform}{zonas de fratura}, os dois lados pertencem à mesma placa e se movem juntos.
\end{solution}

\begin{solution}{exo:b1:rifts-ridges-margins:9}
A taxa de metade é \qty{10.6}{mm.yr^{-1}}, \qty{10.6}{km.Myr^{-1}}. (a) O degrau de idade é
$350/10.6 = \qty{33.0}{Myr}$. (b) No fim da transformante, o fundo de \qty{33.0}{Myr}
($2500 + 350\sqrt{33.0} = \qty{4511}{m}$) encontra fundo de idade 0 (\qty{2500}{m}): um degrau de cerca de \qty{2010}{m}. A \qty{100}{km} além,
as idades são $450/10.6 = 42.5$ e $100/10.6 = \qty{9.4}{Myr}$, as profundidades
4781 e \qty{3575}{m}: um degrau de cerca de \qty{1210}{m}. A \qty{1000}{km} além,
as idades são 127.4 e \qty{94.3}{Myr}, as profundidades 6450 e \qty{5899}{m}:
\qty{550}{m}. (Essas idades passam de \qty{80}{Myr}, onde a lei da raiz quadrada
vai fundo demais; a lei da placa daria um degrau menor.)
\end{solution}

\begin{solution}{exo:b1:rifts-ridges-margins:10}
(a) Com áreas iguais para cada idade, a profundidade média é a média de $d(t)$ em
$[0, A]$:
\[
  \frac1A\int_0^A\bigl(d_r + C\sqrt t\bigr)\mathrm dt = d_r + \frac CA\cdot\frac23A^{3/2}
  = d_r + \frac23C\sqrt A.
\]
(b) Para $A = \qty{180}{Myr}$: $2500 + \frac23\times350\times13.416 =
\qty{5630}{m}$; para \qty{120}{Myr}: $2500 + \frac23\times350\times10.954 =
\qty{5056}{m}$, \qty{574}{m} mais raso. (c) Seja a subida do mar $r$ e o fundo
carregado afunda $s = r\rho_w/(\rho_m - \rho_w)$. A água, inalterada em volume, deve
encher a mesma profundidade sobre a mesma área, logo $r + s = \qty{574}{m}$,
$r\,\rho_m/(\rho_m - \rho_w) = 574$, e $r = 574\times2300/3300 =
\qty{400}{m}$: centenas de metros, só pela taxa de espalhamento.
\end{solution}

\begin{solution}{exo:b1:rifts-ridges-margins:11}
(a) $\sqrt{x - a} = \sqrt x\sqrt{1 - a/x} \approx \sqrt x\,(1 - a/2x)$ para
$a \ll x$, logo
$\Delta d \approx C\sqrt{x/u}\cdot a/(2x) = Ca/(2\sqrt{ux})$.
(b) Em $x = \qty{450}{km}$: exato \qty{1205}{m}, aproximado
$350\times350/(2\sqrt{10.6\times450}) = \qty{887}{m}$; em \qty{1350}{km}:
\qty{550}{m} e \qty{512}{m}; em \qty{3350}{km}: \qty{334}{m} e
\qty{325}{m}. A razão entre o degrau aproximado e o exato depende apenas de
$r = a/x$: $(r/2)/(1 - \sqrt{1 - r})$, que vale 0.9 para $r = 0.36$. A
aproximação é boa até \qty{10}{\percent} além de $x = 350/0.36 \approx
\qty{970}{km}$, cerca de três deslocamentos da dorsal.
\end{solution}

\begin{solution}{exo:b1:rifts-ridges-margins:12}
(a) $d_2 - d_1 = 3671.1 - 3606.8 = \qty{64.3}{m}$ a \qty{10}{Myr};
$4072.8 - 4065.3 = \qty{7.5}{m}$ a 20; $4415.4 - 4417.0 = \qty{-1.7}{m}$ a
30; $4707.5 - 4713.6 = \qty{-6.1}{m}$ a 40. (b) $d_2 - d_1$ é positivo a 20
e negativo a 30, e ambas as leis são contínuas: elas são iguais entre os dois.
Bissecção: a 25, $d_2 - d_1 = +0.9$; a 27.5, $-0.7$; a 26.25, $+0.01$; a
26.875, $-0.3$; a 26.56, $-0.2$; a 26.41, $-0.08$; a 26.33, $-0.04$: a raiz
está entre 26.25 e 26.33, em \qty{26.3}{Myr}, onde ambas dão
\qty{4294}{m}. (c) $d_1' = 175/\sqrt{26.3} = \qty{34.1}{m.Myr^{-1}}$ e
$d_2' = (3200/62.8)\,e^{-26.3/62.8} = \qty{33.5}{m.Myr^{-1}}$. As leis têm
quase o mesmo valor e a mesma inclinação ali, de modo que a curva unida não tem
quebra visível.
\end{solution}

\begin{solution}{pb:b1:rifts-ridges-margins:1}
\textbf{1.} Afastamento $= $ \omterm{def:b1:cross-sections:throw}{rejeito}$/\tan$(\omterm{def:b1:geological-maps:attitude}{mergulho}): F1 \qty{346}{m}, F2 \qty{770}{m},
F3 \qty{2014}{m}, F4 \qty{163}{m}, F5 \qty{462}{m}, F6 \qty{1172}{m}. F3, com
o maior \omterm{def:b1:cross-sections:throw}{rejeito} e o maior \omterm{def:b1:cross-sections:throw}{afastamento horizontal}, é a falha de borda.

\textbf{2.} A extensão é a soma, \qty{4927}{m}, cerca de \qty{4.9}{km}. O
rifte tinha $60 - 4.9 = \qty{55.1}{km}$ de largura antes do falhamento.

\textbf{3.} A normal a um rifte orientado a \ang{020} aponta para \ang{110}; o
movimento, em direção a \ang{105}, faz \ang{5} com ela. A componente perpendicular
é $6.0\cos\ang{5} = \qty{5.98}{mm.yr^{-1}}$.

\textbf{4.} $\qty{4.927}{km}/\qty{5.98}{km.Myr^{-1}} = \qty{0.82}{Myr}$.

\textbf{5.} Em \qty{3.2}{Myr} as placas se afastaram de
$5.98\times3.2 = \qty{19.1}{km}$ através do rifte. As falhas transportam
$4.93/19.1 = \qty{26}{\percent}$ disso.

\textbf{6.} Os outros três quartos, \qty{14.2}{km}, não estão nas falhas da
seção. Eles podem ser absorvidos por diques intrudidos ao longo do eixo, que abrem
o rifte sem falhá-lo, como em Afar, ou por falhas fora da seção.
Com diques de \qty{8}{m} de largura, $\qty{14200}{m}/\qty{8}{m} \approx 1780$ diques são
necessários, um a cada $3.2\times10^6/1780 \approx 1800$ anos em média.

\textbf{7.} De C2An: $(21.0 + 20.2)/2.581 = 41.2/2.581 =
\qty{15.96}{mm.yr^{-1}}$. De C3n: $(41.0 + 40.2)/5.235 = 81.2/5.235 =
\qty{15.51}{mm.yr^{-1}}$.

\textbf{8.} $(41.2 - 2.2)/2.581 = \qty{15.11}{mm.yr^{-1}}$ e
$(81.2 - 2.2)/5.235 = \qty{15.09}{mm.yr^{-1}}$: corrigidas, as duas concordam a
\qty{0.02}{mm.yr^{-1}}. A correção importa mais para o cron mais jovem, cuja
distância é menor.

\textbf{9.} O deslocamento para fora desloca as duas bordas de um flanco na mesma
medida, logo ele se cancela na diferença. Somando os dois flancos,
$(81.2 - 41.2)/(5.235 - 2.581) = 40.0/2.654 = \qty{15.07}{mm.yr^{-1}}$.

\textbf{10.} Nordeste: $(41.0 - 21.0)/2.654 = \qty{7.54}{mm.yr^{-1}}$;
sudoeste: $(40.2 - 20.2)/2.654 = \qty{7.54}{mm.yr^{-1}}$. O espalhamento é
simétrico; a diferença de \qty{0.8}{km} entre os flancos, a mesma para os dois
crons, é um deslocamento constante, como um eixo colocado \qty{0.4}{km} fora
do centro.

\textbf{11.} $\cos\Delta = \sin\ang{30.9}\sin\ang{18.0} +
\cos\ang{30.9}\cos\ang{18.0}\cos\ang{16.4} = 0.15870 + 0.78287 = 0.94157$, logo
$\Delta = \ang{19.69}$. Com $\omega R = 0.403\times(\pi/180)\times6371 =
\qty{44.81}{mm.yr^{-1}}$, a taxa é $44.81\sin\ang{19.69} =
\qty{15.10}{mm.yr^{-1}}$, as taxas corrigidas das questões 8 e 9.

\textbf{12.} $\cos\Delta = \sin\ang{30.9}\sin\ang{26.0} +
\cos\ang{30.9}\cos\ang{26.0}\cos\ang{11.4} = 0.22512 + 0.75600 = 0.98112$, logo
$\Delta = \ang{11.15}$ e a taxa é $44.81\sin\ang{11.15} =
\qty{8.7}{mm.yr^{-1}}$.

\textbf{13.} As previsões, 15.1 e \qty{8.7}{mm.yr^{-1}}, concordam com as
medidas 15 e \qty{9}{mm.yr^{-1}}. A taxa é $\omega R\sin\Delta$, e o
polo fica a noroeste do mar Vermelho, no nordeste da Líbia: quanto mais ao sul, mais longe
do polo, mais rápida a abertura.

\textbf{14.} $15.09\times5.235 = \qty{79}{km}$ de fundo oceânico.

\textbf{15.} $5000/120.6 = \qty{41}{km.Myr^{-1}}$, \qty{41}{mm.yr^{-1}} em
média.

\textbf{16.} $5000/15.1 = \qty{331}{Myr}$ à taxa central;
$5000/8.7 = \qty{575}{Myr}$ à taxa do norte.

\textbf{17.} Abrindo-se em torno de um polo próximo, o mar Vermelho alarga-se como um par de
tesouras: sua largura cresce em proporção a $\sin\Delta$, e ele continua uma cunha
que se estreita em direção ao polo. A \qty{8.7}{mm.yr^{-1}}, a parte norte precisa de
$79/8.7 = \qty{9.1}{Myr}$ para formar os \qty{79}{km} de fundo oceânico que a
parte central tem agora.

\textbf{18.} A \qty{15}{mm.yr^{-1}}, o eixo do mar Vermelho é uma dorsal ultralenta; um
vale axial entre falhas normais é o esperado.

\textbf{19.} $2500 + 350\times1 = \qty{2850}{m}$;
$2500 + 350\sqrt{10} = \qty{3607}{m}$; $2500 + 350\sqrt{40} = \qty{4714}{m}$.

\textbf{20.} $t = (1500/350)^2 = \qty{18.4}{Myr}$.

\textbf{21.} Nos primeiros \qty{10}{Myr} ele afunda $350\sqrt{10} = \qty{1107}{m}$;
entre 70 e \qty{80}{Myr}, $350(\sqrt{80} - \sqrt{70}) = 350\times(8.944 -
8.367) = \qty{202}{m}$, cinco vezes menos no mesmo intervalo de tempo.

\textbf{22.} $\mathrm dd/\mathrm dt = C/(2\sqrt t) = 175/\sqrt{80} =
\qty{19.6}{m.Myr^{-1}}$, cerca de \qty{0.02}{mm} por ano.

\textbf{23.} $490/\sqrt{80} = \qty{54.8}{mW.m^{-2}}$.

\textbf{24.} Lei da raiz quadrada: $2500 + 350\sqrt{80} = 2500 + 350\times8.944 =
\qty{5630}{m}$. Lei da placa: $6400 - 3200\,e^{-80/62.8} = 6400 - 3200\times0.2798
= \qty{5505}{m}$. Elas diferem em \qty{126}{m}. A \qty{80}{Myr} as duas leis ainda
valem: a lei da raiz quadrada até cerca de \qty{80}{Myr}, a lei da placa a partir de cerca de
\qty{20}{Myr}; além disso, a lei da raiz quadrada vai fundo demais.

\textbf{25.} A lei da raiz quadrada dá \qty{5630}{m}: com aproximação de um décimo de
quilômetro, o fundo do novo oceano fica a uma profundidade de
\textbf{\qty[text-series-to-math=true]{5.6}{km}} quando tem 80 milhões de anos
(a lei da placa dá \qty{5.5}{km}).
\end{solution}
